\section{Introducción}

El transformador monofásico constituye una de las máquinas eléctricas estáticas de mayor relevancia en los sistemas de transmisión, distribución y utilización de la energía eléctrica. Su estudio experimental y modelado circuital se ha desarrollado progresivamente en las prácticas previas: en la Práctica Nº 4 se determinaron los parámetros de su circuito equivalente (relación de transformación $a$, resistencia y reactancia de dispersión, y la rama de excitación compuesta por la resistencia de pérdidas en el núcleo $R_c$ y la reactancia de magnetización $X_m$), mientras que en la Práctica Nº 5 se evaluó su comportamiento y regulación de tensión bajo distintas condiciones de carga.

El presente informe se centra en la determinación y el análisis del \textbf{rendimiento} (o eficiencia) de dicha unidad bajo condiciones reales de operación en el laboratorio. El rendimiento cuantifica la efectividad con la cual el transformador transfiere la potencia activa absorbida desde la red de alimentación hacia la carga conectada en sus bornes secundarios. Formalmente, se define como la razón entre la potencia activa de salida y la potencia activa de entrada:

\begin{equation}
    \eta = \frac{P_{\text{salida}}}{P_{\text{entrada}}} = \frac{P_{\text{salida}}}{P_{\text{salida}} + P_{\text{perd}}}
\end{equation}

El análisis experimental del rendimiento reviste un interés fundamental en la ingeniería eléctrica, pues no se limita al cálculo de una cifra porcentual de eficiencia, sino que permite comprender la distribución de las pérdidas internas de la máquina. La energía disipada se descompone rigurosamente en dos grupos de naturalezas físicas distintas: las \textbf{pérdidas sin carga o en el hierro} ($P_{\text{Fe}}$), asociadas a los fenómenos de histéresis magnética y corrientes parásitas de Foucault en el núcleo, y las \textbf{pérdidas con carga o en el cobre} ($P_{\text{Cu}} = I_2^2 R_{\text{eq2}}$), causadas por el efecto Joule en los arrollamientos primario y secundario.

A través de esta separación de pérdidas, se evalúa experimentalmente la variación de la eficiencia frente a distintos regímenes de carga (resistiva e inductiva), identificando el \textbf{punto de rendimiento máximo}, en el cual las pérdidas en el hierro se igualan a las pérdidas en el cobre. Dicho punto permite justificar criterios de diseño, dimensionamiento óptimo y estimación de la potencia nominal de la unidad de prueba. Asimismo, se incorpora el método de la resistencia para registrar la elevación de temperatura del bobinado tras el ensayo a plena carga y efectuar las debidas correcciones térmicas a la resistencia de los devanados, siguiendo las prescripciones sobre aumento de temperatura de la norma venezolana COVENIN 3172-95 (punto 5.7) y COVENIN 536-94, así como la normativa internacional ANSI/IEEE C57.12.91-2001.
